% !TeX program = xelatex
\documentclass[11pt]{article}
\usepackage{acl}
\usepackage{times}
\usepackage{latexsym}
\usepackage[T1]{fontenc}
\usepackage{microtype}
\usepackage{inconsolata}
\usepackage{xcolor}
\usepackage{graphicx}
\graphicspath{{images/}}
\usepackage{url}
\usepackage{booktabs}
\usepackage{array}
\usepackage{multirow}
\usepackage{mwe}
\usepackage{ifxetex}
\ifxetex
  \usepackage{fontspec}
  \newfontfamily\devanagarifont[Script=Devanagari]{FreeSerif}
  \newcommand{\skt}[1]{{\devanagarifont #1}}
\else
  \usepackage[utf8]{inputenc}
  \newcommand{\skt}[1]{#1}
\fi
\setlength{\tabcolsep}{4pt}
\usepackage{makecell}
\renewcommand{\theadfont}{\small\bfseries}

\title{SanskritRasaBank: Benchmarking Multilingual Transformers on
Fine-Grained Affective Classification in Epic and Vedic Corpora}

\author{Suhrit Ghimire\textsuperscript{1}  \quad Minni Jain\textsuperscript{1} 
\quad Rohini Raj Timilsina\textsuperscript{2}\\
  \textsuperscript{1}Delhi Technological University, India \\
  \textsuperscript{2}Department of Sanskrit, Tribhuvan University, Nepal \\
  \texttt{suhrit\_23ec245@dtu.ac.in, minnijain@dtu.ac.in} \\ \url{}}

\begin{document}
\maketitle

\begin{abstract}
Bharata Muni's Nāṭyaśāstra first codified the nine rasas, aesthetic essences ranging from love (Śṛṇgāra) to serenity (Śānta), that define the traditional Indian literary experience. While this framework has anchored Sanskrit literary criticism for over two millennia, it has remained largely inaccessible to computational linguistics. The absence of any large-scale annotated corpus has left the emotional architecture of Sanskrit's greatest texts empirically unexplored.
To bridge this gap, we introduce SanskritRasaBank, the first large-scale, expert-validated corpus for nine-class rasa classification. The dataset comprises 17,462 meticulously verified verse annotations drawn from eight classical sources. Our annotation pipeline began with an initial labeling pass by a three-LLM ensemble (GPT-4o, DeepSeek-Chat, and LLaMA-3.1-8b-instant) governed by strict consensus rules, followed by comprehensive human validation by a trained Sanskrit philologist and his team, producing a high-precision silver standard.
We benchmark classical machine learning (SVM with character n-gram TF-IDF) alongside four transformer architectures across two experimental phases. In baseline 3-fold cross-validation (Phase 1), MuRIL achieves 80.48% accuracy and 80.60% weighted F1, significantly outperforming IndicBERT V2 (65.21%) and XLM-RoBERTa (59.28%). Following targeted hyperparameter optimization (Phase 2), IndicBERT V2 surges to 81.45% while MuRIL stabilizes at 80.65%.
Applying our best-performing model, we conduct the first large-scale computational emotional mapping of the complete Vālmīki Rāmāyaṇa (23,402 verses across all seven kāṇḍas) and Mandala 1 of the Ṛgveda (1,758 hymns), revealing distinct kanda-wise affective arcs and genre-specific rasa signatures that empirically validate centuries of traditional literary scholarship. The corpus focuses entirely on padhyas (metrical verses). All data and models are publicly available at \url{https://github.com/SuhritGhimire/SanskritRasaBank}.
\end{abstract}

\vspace{0.5em}
\noindent \textbf{Keywords:} SanskritRasaBank, Nava-Rasa, Sanskrit NLP, Affective Computing, Multi-Label Classification, LLM-Ensemble, Expert-Validated Corpus, Digital Humanities, Indic Transformers, Low-Resource NLP.

\section{Introduction}

Language encodes more than mere propositions; it conveys a distinct aesthetic taste. Ancient Indian literary theorists termed this relished quality \textit{rasa}, defining it as the emotional flavor of a literary moment and the felt essence of a verse that transcends its referential content. In the \textit{N\=a\d tyaś\=astra} (c.~200 BCE), Bharata Muni codified nine such flavors: love (\textit{Ś\d r\d ng\=ara}), laughter (\textit{H\=asya}), grief (\textit{Karu\\d n\=a}), fury (\textit{Raudra}), heroism (\textit{V\=ira}), terror (\textit{Bhay\=anaka}), disgust (\textit{B\=ibhatsa}), wonder (\textit{Adbhuta}), and serenity (\textit{Ś\=anta}). While this framework has served as the foundational language of Sanskrit criticism for millennia, the concept of rasa has been largely absent from the modern computational record. There are no sizable annotated corpora available, meaning the underlying emotional structures of Sanskrit's greatest texts have never been mapped at scale.

This paper bridges that disciplinary divide. We begin by introducing \texttt{SanskritRasaBank}, a newly constructed corpus of 17,462 expert-verified verses. To build this dataset, we developed a hybrid pipeline where an ensemble of three large language models (GPT-4o, DeepSeek-Chat, LLaMA-3.1-8b-instant) provided initial labels under strict consensus filtering protocols. These labels were then subjected to a meticulous human audit by experts in Sanskrit poetics. Leveraging this new resource, we provide a thorough benchmark of classical and neural classification approaches. Moving from character n-gram SVM baselines to specialized multilingual transformers, we demonstrate that MuRIL's Indian-language pre-training gives it a decisive edge on this morphologically dense and culturally nuanced task. Finally, we put our best-performing model to work, performing the first computational emotional cartography of the complete \textit{V\=alm\=iki R\=am\=aya\d na} and the first mandala of the \textit{\d Rgveda}. The resulting affective arcs not only validate but significantly enrich the traditional scholarly reception of these canonical texts.

Specifically, our findings reveal that while \textit{V\=ira} dominates six of the seven k\=a\d n\d das of the R\=am\=aya\d na, the Ayodhy\=a K\=a\d n\d da is uniquely and overwhelmingly \textit{Karu\d n\=a}-dominant (33.3\%), standing in precise accord with its traditional epithet \textit{karu\d na-pradh\=ana}. The Yuddha K\=a\d n\d da peaks at 39.0\% \textit{V\=ira} with a substantial \textit{Raudra} undertow. The \textit{\d Rgveda}, by contrast, presents an entirely different emotional grammar. Here, \textit{Adbhuta} (Wonder) is dominant at 23.78\%, while \textit{Ś\d r\d ng\=ara} (15.13\%) and \textit{Ś\=anta} (12.12\%) are both elevated relative to the epic, capturing the hymnal aesthetic of divine address and cosmic marveling.

\section{Related Work}

\subsection{Sanskrit Natural Language Processing}

Sanskrit NLP finds its computational roots in Huet and collaborators' Sanskrit Heritage Platform (Huet, 2003; Huet and Kulkarni, 2009), a morphological analyzer covering thousands of roots and their inflectional paradigms. The subsequent decade produced foundational tools: \textit{sandhi} splitters (Mittal, 2010; Hellwig, 2015), dependency parsers (Kulkarni et al., 2010), and statistical word segmenters (Hellwig and Nehrdich, 2018; Krishna et al., 2018a). The Digital Corpus of Sanskrit (DCS; Hellwig, 2010 onward) now aggregates over five million words. Neural methods arrived with Sanskrit word embeddings (Satuluri and Bhatt, 2020), followed by transformer fine-tuning (Krishna et al., 2021). SanBERTa (Pal et al., 2023) is the first Sanskrit-native transformer. On the affective side, the IIT-R STSA dataset (Kumar et al., 2022) provides 12,892 R\=am\=aya\d na sentences with binary/ternary sentiment, not nine-class rasa labels. Sandhan et al.\ (2025) recently proposed an interpretable framework for Sanskrit poetic elements including mood, yet large-scale rasa corpora remain absent.

\subsection{Rasa Theory and Computational Aesthetics}

Our theoretical foundation rests on Bharata Muni's \textit{N\=a\d tyaś\=astra} (c.\ 200 BCE), which articulates aesthetic experience through a tripartite mechanism: \textit{vibh\=ava} (the objective stimulus), \textit{anubh\=ava} (the outward expression), and \textit{vyabhic\=aribh\=ava} (transitory affective states). In the tenth century, Abhinavagupta's \textit{Abhinavabh\=arat\=i} expanded this framework into a comprehensive phenomenology of \textit{ras\=asv\=ada}, the active relishing of art. These classical principles form the conceptual scaffolding that our annotation protocol operationalizes. While work on rasa classification in neighboring Indic languages offers valuable precedents, such as the \textit{Bhaavana} system's feature-based classification of Hindi poetry (Mishra et al., 2018), SVM-based assignments for Hindi verse (Kumar et al., 2019), and emotion data from Hindi fiction in the BHAAV corpus (Singla et al., 2019), Sanskrit presents unique challenges. Its morphophonological complexity, distinct aesthetic register, and the sheer scale of texts like the 24,000-verse R\=am\=aya\d na demand a dedicated, language-specific approach. Furthermore, our annotation design is guided by the ethical and methodological frameworks for automatic emotion recognition proposed by Mohammad (2022).

\subsection{LLM-Ensemble Annotation}

The use of large language models for dataset annotation has rapidly transitioned from a novel experiment to an established methodology. Wang et al.\ (2023) demonstrated that LLM-generated data can successfully bootstrap classifiers, while others have found LLMs capable of matching or even outperforming human crowd-workers on standard annotation tasks (Honovich et al., 2023; Gilardi et al., 2023). However, translating these successes to \emph{culturally embedded, highly specialized domains} is far from automatic (M{\o}ller et al., 2023; Qin et al., 2023). Rasa classification represents exactly this kind of specialized task, demanding a nuanced internalization of Sanskrit poetics, culturally specific emotional registers, and a 2,000-year-old aesthetic tradition that general-purpose LLMs may only capture superficially. To mitigate these risks, we employ a stringent two-layer validation architecture. First, we require \textbf{unanimous consensus across a three-LLM ensemble}, treating any single model's prediction as insufficient on its own. Second, we mandate \textbf{expert human validation} conducted by a trained Sanskrit philologist and his graduate team. Together, these paired mechanisms elevate LLM annotation from a convenient approximation to a robust, independently audited methodology.

\subsection{Multilingual and Indic NLP}

mBERT (Devlin et al., 2019) established cross-lingual transfer for 104 languages. XLM and XLM-RoBERTa (Lample and Conneau, 2019; Conneau et al., 2020) strengthened multilingual representations through masked language modeling at 100-language scale. The Indic-specific line of work, encompassing IndicBERT (Kakwani et al., 2020) and IndicBERTV2 (Doddapaneni et al., 2023), focused pre-training on 12 Indian languages. MuRIL (Khanuja et al., 2021) goes further: a BERT-large architecture trained on 17 Indian languages \emph{and} their transliterated variants, making it the only model in this landscape that explicitly includes Sanskrit in its pre-training corpus. Recent work by Sandhan et al.\ (2022) on subword tokenization for Sanskrit further establishes the viability of transformer pre-training for classical Indic language tasks.

\section{Dataset Construction}

\subsection{Source Material}

The corpus is assembled from eight classical Sanskrit sources to ensure broad rasa coverage. The primary source is the \textit{V\=alm\=iki R\=am\=aya\d na} (all seven k\=a\d n\d das, approximately 23,400 verses), which provides strong coverage of \textit{V\=ira}, \textit{Karu\d n\=a}, and \textit{Raudra} but is structurally sparse on \textit{H\=asya} and \textit{B\=ibhatsa}. To redress this imbalance, we supplemented with the following texts:

\begin{itemize}
    \item \textit{Kath\=asarits\=agara} (Ocean of the Streams of Story): rich in narrative variety, contributing \textit{Adbhuta} and \textit{Bhay\=anaka} registers
    \item \textit{Pa\~ncatantra}: didactic fables supplying \textit{H\=asya}, irony, and \textit{B\=ibhatsa}
    \item \textit{Amaru\'{s}ataka}: erotic lyric poetry, the densest source of \textit{Ś\d r\d ng\=ara}
    \item \textit{Vet\=alapañcavi\d m\'{s}ati}: supernatural narratives rich in \textit{Bhay\=anaka} and the grotesque
    \item \textit{Abhij\~n\=anaś\=akuntalam} (K\=alid\=asa): classical drama with layered \textit{Ś\d r\d ng\=ara} and \textit{Karu\d n\=a}
    \item Mandala 1 of the \textit{\d Rgveda}: Vedic hymns providing hymnal \textit{Adbhuta} and \textit{Ś\=anta}
    \item Selected narrative excerpts from GRETIL digital repositories
\end{itemize}
Our data collection is restricted to \textit{padhyas} (verses), as the rhythmic and structural constraints of Sanskrit poetry create a distinct affective environment compared to \textit{gadhyas} (prose).

The initial pool comprised approximately 29,900 verses before consensus filtering.

\subsection{Annotation Pipeline}

We implemented a three-stage pipeline (visualised in Figure~\ref{fig:pipeline}) designed to maximise precision while managing the inherent ambiguity of rasa classification.

\textbf{Stage 1: LLM Ensemble Labeling.} Each verse was independently submitted to three large language models: \textbf{GPT-4o} (OpenAI, 2024), \textbf{DeepSeek-Chat} (DeepSeek-AI, 2024), and \textbf{LLaMA-3.1-8b-instant} (Meta AI, 2024). Each model was prompted via a structured instruction to identify the single dominant rasa and supply a focus keyword grounding the prediction in textual evidence. All labeling was performed independently with no inter-model communication.

\textbf{Stage 2: Consensus Filtering.} After the initial labeling pass, we applied rasa-specific consensus rules:
\begin{itemize}
    \item \textbf{Majority Rasas} (\textit{V\=ira}, \textit{Karu\d n\=a}, \textit{Adbhuta}, \textit{Raudra}, \textit{Ś\d r\d ng\=ara}, \textit{Ś\=anta}): Only verses where all three LLMs assigned the same label were retained. Unanimous three-way agreement serves as strong evidence that the rasa signal is unambiguous and cross-model consistent.
    \item \textbf{Minority Rasas} (\textit{H\=asya}, \textit{B\=ibhatsa}, \textit{Bhay\=anaka}): For structurally rare classes, verses where at least two of three LLMs agreed were retained, preventing near-total exclusion of these categories and maintaining training viability.
\end{itemize}
This dual-threshold consensus design balances dataset precision with class coverage.

\textbf{Stage 3: Expert Human Validation.} The consensus-filtered dataset underwent a rigorous audit by Rohini Raj Timilsina (Lecturer, Department of Sanskrit, Tribhuvan University) and his team of graduate students, all formally trained in classical Sanskrit poetics. Each label was individually evaluated against its source verse. Annotations found to be contextually inaccurate, ambiguous, or misaligned with traditional rasa theory were systematically discarded. This painstaking human review refined the dataset into the final \texttt{MERGED\_FINAL} corpus of \textbf{17,462 verified samples}, yielding what is currently the largest and most rigorously validated Sanskrit rasa resource available.

\begin{figure*}[!t]
\centering
\includegraphics[width=\textwidth]{images/image.png}
\caption{Three-stage annotation pipeline for \texttt{SanskritRasaBank}: LLM-ensemble labeling, multi-threshold consensus filtering, and expert human validation by Sanskrit philologists.}
\label{fig:pipeline}
\end{figure*}

The efficacy of this multi-stage annotation pipeline (Figure~\ref{fig:pipeline}) is directly reflected in the resulting dataset's structural integrity. By enforcing unanimous consensus for the majority classes and expert-audited high-precision labels throughout, we arrive at a \texttt{MERGED\_FINAL} distribution (Figure~\ref{fig:class_dist}) where class boundaries are sharp and philologically consistent. This rigorous methodological foundation ensures that the dataset captures the authentic emotional architecture of the source texts, providing a stable and reliable ground truth for the comprehensive model benchmarking which follows.

\subsection{Multi-Rasa Architecture}

Unlike prior single-label classification approaches, our inference pipeline captures the top-3 predicted rasas with confidence scores for each verse. The primary prediction (\textit{rasa\_1}) represents the dominant rasa and is used for all corpus-level analyses. The secondary (\textit{rasa\_2}) and tertiary (\textit{rasa\_3}) predictions are provided to capture \textit{rasa-dhvani}, the resonant undercurrents that classical critics identify as essential to the full aesthetic experience of a verse. These secondary predictions reveal systematic co-occurrence patterns in the corpus (for example, the near-universal \textit{V\=ira-H\=asya} secondary pairing in battle verse, or the \textit{Karu\d n\=a-Ś\=anta} blend in bereavement passages) that a single-label approach would erase.

\subsection{Final Corpus Distribution}

The verified dataset comprises 17,462 data points. Table~\ref{tab:final_dist} shows the final distribution.

\begin{table}[htbp]
\centering
\small
\caption{Verified Dataset Class Distribution (\texttt{MERGED\_FINAL}).}
\label{tab:final_dist}
\begin{tabular}{lrc}
\toprule
Rasa & Count & \% \\
\midrule
V\=ira       & 6,475 & 37.1\% \\
Karu\d n\=a  & 3,948 & 22.6\% \\
Adbhuta      & 3,159 & 18.1\% \\
Raudra       & 2,999 & 17.2\% \\
Ś\d r\d ng\=ara & 2,672 & 15.3\% \\
Ś\=anta      & 2,087 & 12.0\% \\
Bhay\=anaka & 1,105 &  6.3\% \\
B\=ibhatsa   &   532 &  3.0\% \\
H\=asya      &   425 &  2.4\% \\
\midrule
\textbf{Total} & \textbf{17,462} & \textbf{100\%} \\
\bottomrule
\end{tabular}
\end{table}

\begin{figure*}[!t]
\centering
\includegraphics[width=\textwidth]{images/Unknown.png}
\caption{Distribution of rasa labels in \texttt{MERGED\_FINAL}. \textit{V\=ira} and \textit{Karu\d n\=a} are the most frequent classes, reflecting the R\=am\=aya\d na as primary source; minority classes (\textit{H\=asya}, \textit{B\=ibhatsa}) were boosted via supplementary texts.}
\label{fig:class_dist}
\end{figure*}

\section{Models and Training}

\subsection{Classical Baseline: SVM with Character N-gram TF-IDF}

To establish a rigorous classical baseline, we trained a Linear Support Vector Machine (LinearSVC) using a character n-gram TF-IDF representation (spanning 2- to 5-grams, capped at a 200,000 vocabulary size, and applying sublinear TF scaling) with probabilities calibrated via Platt scaling. Character n-grams are uniquely suited to Sanskrit; they implicitly capture \textit{sandhi} combinations and morphological sub-word patterns without demanding explicit word segmentation. We addressed class imbalance by applying balanced class weighting, and evaluated the model using 5-fold stratified cross-validation on the \texttt{MERGED\_FINAL} dataset. While computationally efficient, the SVM performance isolates a fundamental limitation: surface-level features capture morphological regularities but fail to synthesize the deep context and suprasegmental dependencies required to interpret classical rasa.

\subsection{Transformer Models}

We evaluate four transformer architectures:

We evaluate four distinct transformer architectures:

\textbf{MuRIL (google/muril-large-cased)}: A 24-layer BERT-large architecture (236M parameters) pre-trained by Google on 17 Indian languages and their transliterated scripts. Notably, Sanskrit is explicitly included in this pre-training corpus, rendering MuRIL the only model in our benchmark with direct, pre-trained exposure to Sanskrit text.

\textbf{IndicBERT V2 (ai4bharat/IndicBERTv2-MLM-Sam-TLM)}: A 12-layer RoBERTa-style model pre-trained on a massive 35-billion-token corpus spanning 12 Indian languages.

\textbf{XLM-RoBERTa-large (xlm-roberta-large)}: A comprehensive multilingual RoBERTa model (24 layers, 560M parameters) pre-trained across 100 languages using masked language modeling.

\textbf{SanBERTa}: A Sanskrit-native RoBERTa model pre-trained exclusively on Sanskrit corpora. We include it to test the hypothesis that monolingual specialization outweighs broad multilingual exposure.

\subsection{Training Configuration}

\textbf{Phase 1 Baseline 3-fold Cross-Validation.} In the first phase, all models were trained under their default recommended hyperparameters to establish baseline performance: learning rate 2e-5, batch size 32, up to 15 epochs, maximum sequence length 256, Focal Loss with class weights. MuRIL used full fine-tuning; IndicBERT V2 and XLM-RoBERTa used LoRA (Hu et al., 2022) for parameter efficiency.

\textbf{Phase 2 Hyperparameter Optimisation (Single Fold).} Following Phase 1, where IndicBERT and XLM-RoBERTa showed substantially sub-optimal performance, we conducted a targeted hyperparameter search on a single representative fold to maximise resource efficiency. For MuRIL, we implemented:
\begin{itemize}
    \item Focal Loss (Lin et al., 2017; FL($\gamma$=2)) combined with label smoothing ($\epsilon$=0.1) to handle both class imbalance and residual LLM annotation noise
    \item Layer-Wise Learning Rate Decay (LLRD) with a 0.95 decay factor per BERT layer (top layers: 1e-5; embedding layer: $\sim$3.6e-7)
    \item Cosine learning rate schedule with a 15\% warmup ratio
    \item Gradient checkpointing and mixed precision (FP16) for memory efficiency
    \item Early stopping with patience=4 on weighted F1; model saved at peak validation F1
\end{itemize}
For IndicBERT V2 and XLM-RoBERTa in Phase 2, learning rates and LoRA rank parameters were tuned, and Focal Loss was introduced. All models used stratified splitting (random\_state=42) ensuring reproducible class distributions across folds.

\section{Experimental Results}

\subsection{Phase 1: Baseline 3-Fold Cross-Validation}

Table~\ref{tab:phase1_results} reports the mean and standard deviation of accuracy and weighted F1 across three folds.

\begin{table}[htbp]
\centering
\small
\caption{Phase 1: 3-fold CV (Mean $\pm$ Std \%). SVM (TF-IDF) baseline was substantially lower.}
\label{tab:phase1_results}
\begin{tabular}{@{}lccc@{}}
\toprule
Model & Acc & W-F1 & M-F1 \\
\midrule
MuRIL        & \textbf{80.5$\pm$0.5} & \textbf{80.6$\pm$0.5} & \textbf{75.7$\pm$1.2} \\
IndicBERTv2  & 65.2$\pm$1.0          & 66.5$\pm$1.0          & 56.8$\pm$1.1          \\
XLM-R        & 59.3$\pm$0.7          & 60.8$\pm$1.2          & 51.8$\pm$0.9          \\
SanBERTa     & 58.1$\pm$1.0          & 59.3$\pm$0.9          & 49.8$\pm$1.1          \\
\bottomrule
\end{tabular}
\end{table}

MuRIL's commanding lead in Phase 1 outpacing IndicBERT V2 by over 15 percentage points and XLM-RoBERTa by over 21 points stems directly from its Sanskrit-inclusive pre-training. It appears that standard multilingual models, irrespective of their massive parameter counts, lack adequate exposure to the specific phonological and morphological structures of classical Sanskrit. Interestingly, SanBERTa underperforms all the multilingual models. This counters the intuitive assumption that a Sanskrit-native model should naturally excel, and instead suggests that the immense scale and linguistic diversity of modern multilingual corpora provide a representational advantage that localized, monolingual training cannot easily match.

Meanwhile, the SVM baseline, although uncompetitive against transformer architectures on a complex nine-class task, serves a critical diagnostic purpose. Its character n-gram features successfully recognize surface morphological markers (such as k\=ara\d na forms, sandhi junctions, and case suffixes), yet wholly fail to grasp the contextual depth requisite for rasa interpretation. Rasa transcends the lexicon; it is a holistic property born from imagery, rhythm, and narrative positioning. Transformers have the capacity to model this distributed context, precisely where surface-level algorithms fall short.

\begin{table}[htbp]
\centering
\small
\caption{Phase 1: Per-class F1 for MuRIL (3-fold average).}
\label{tab:muril_perclass_p1}
\begin{tabular}{lccc}
\toprule
Rasa & Precision & Recall & F1 \\
\midrule
Ś\d r\d ng\=ara & 0.886 & 0.877 & 0.882 \\
H\=asya         & 0.557 & 0.670 & 0.607 \\
Karu\d n\=a     & 0.853 & 0.810 & 0.831 \\
Raudra          & 0.773 & 0.753 & 0.763 \\
V\=ira          & 0.833 & 0.808 & 0.820 \\
Bhay\=anaka   & 0.583 & 0.720 & 0.636 \\
B\=ibhatsa      & 0.620 & 0.678 & 0.647 \\
Adbhuta         & 0.827 & 0.798 & 0.812 \\
Ś\=anta         & 0.689 & 0.847 & 0.756 \\
\midrule
\textbf{Weighted avg} & \textbf{0.807} & \textbf{0.805} & \textbf{0.806} \\
\bottomrule
\end{tabular}
\end{table}

Noteworthy per-class patterns in Phase 1: \textit{Ś\d r\d ng\=ara} is consistently the best-classified rasa (F1=0.882), owing to its distinctive lexical and imagistic vocabulary. \textit{H\=asya} and \textit{Bhay\=anaka} are the weakest (F1=0.607 and 0.636 respectively), reflecting both their minority class status and the genuine semantic proximity of \textit{Bhay\=anaka} to \textit{Raudra} (both involve violent intensity) and \textit{H\=asya} to \textit{Ś\d r\d ng\=ara} (both involve social playfulness). These confusions are precisely what a human annotator also finds hardest to resolve.

\subsection{Phase 2: Hyperparameter Optimisation}

Table~\ref{tab:phase2_results} reports Phase 2 results on a single held-out fold; the complete performance trajectory across both phases and all models is summarised in Figure~\ref{fig:model_benchmark}.

\begin{table}[htbp]
\centering
\small
\caption{Phase 2: After hyperparameter optimisation.}
\label{tab:phase2_results}
\begin{tabular}{@{}lccc@{}}
\toprule
Model & Acc & W-F1 & M-F1 \\
\midrule
IndicBERTv2 & \textbf{81.5} & \textbf{81.5} & \textbf{76.7} \\
MuRIL       & 80.7          & 80.7          & 76.5          \\
XLM-R       & 78.9          & 78.8          & 73.6          \\
\bottomrule
\end{tabular}
\end{table}

The gains achieved in Phase 2 are remarkably pronounced, illuminating distinct behaviors across the architectures. IndicBERT V2 exhibits the most dramatic leap, surging by \textbf{16.24 percentage points} (from 65.21\% to 81.45\%). This extraordinary improvement suggests that while IndicBERT's underlying Indian-language representations are fundamentally suited to Sanskrit, its default LoRA configuration was too heavily regularized and its standard learning rate misaligned for this specialist domain. By relaxing these constraints through careful tuning, the model swiftly converged to near-optimal performance.

XLM-RoBERTa demonstrates an even steeper raw gain of \textbf{19.59 percentage points} (from 59.28\% to 78.87\%). Although it is the largest model evaluated (560M parameters), its initial struggles highlight the inherent difficulty of adapting a generalized multilingual model to a low-resource classical language without dedicated domain tuning. Once properly calibrated in Phase 2, XLM-RoBERTa closed the gap to within 1.78 percentage points of the leading model.

By contrast, MuRIL's improvement is marginal (\textbf{0.17 pp}, moving from 80.48\% to 80.65\%). This plateau is entirely expected for a model whose pre-training already aligns closely with the target linguistic domain. MuRIL effectively operates near mathematically optimal levels right out of the box. From a tactical standpoint, this robustness is a major asset: MuRIL delivers highly reliable performance without demanding exhaustive, compute-heavy hyperparameter searches, making it the premier choice for practitioners.

\begin{figure}[!ht]
\centering
\includegraphics[width=\columnwidth]{images/Unknown-2.png}
\caption{Accuracy improvement from Phase 1 to Phase 2 for each model. IndicBERT and XLM-RoBERTa show dramatic gains; MuRIL is stable, reflecting its pre-training advantage.}
\label{fig:phase_compare}
\end{figure}

\begin{figure*}[!t]
\centering
\includegraphics[width=\textwidth]{images/Unknown-3.png}
\caption{Per-class F1 scores for MuRIL V2 (sorted high to low). \textit{V\=ira}, \textit{Ś\d r\d ng\=ara}, and \textit{Karu\d n\=a} achieve the highest F1; \textit{H\=asya} and \textit{Bhay\=anaka} remain the most challenging due to their minority status and semantic adjacency to neighbouring rasas. Ladder-style annotations ensure all values are legible; the RdYlGn colour scale encodes performance.}
\label{fig:perclass}
\end{figure*}

While the aggregate performance metrics (Weighted F1) show a clear hierarchy among models, the per-class breakdown (Figure~\ref{fig:perclass}) reveals the underlying source of these gains. The transition from the Phase 1 baselines to the highly optimized Phase 2 models, as summarized in the comprehensive benchmark trajectory (Figure~\ref{fig:model_benchmark}), is driven largely by the models' improved ability to distinguish semantically adjacent minority classes like \textit{H\=asya} and \textit{Bhay\=anaka}. By fine-tuning hyperparameters and introducing class-weighted focal loss, we effectively bridge the performance gap for structurally rare rasas. The "ladder" of F1 scores demonstrates that MuRIL and the optimized IndicBERT V2 have moved beyond surface lexical matching to capture the nuanced affective signals required for reliable Sanskrit literary analysis, providing a high-resolution view of the models' representational depth.

\begin{figure*}[!t]
\centering
\includegraphics[width=\textwidth]{images/Unknown-7.png}
\caption{Full model benchmark across all architectures and both training phases, ranked by weighted F1. MuRIL Phase 1 already exceeds all other Phase 1 baselines; the SVM classical baseline is included for reference. The dashed line marks the overall best result (IndicBERT V2 Phase 2, 80.85\%).}
\label{fig:model_benchmark}
\end{figure*}

\subsection{Model Selection for Corpus Inference}

Based on its unparalleled cross-validated stability (Phase 1 3-fold mean of 80.48\% with a standard deviation of just $\pm$0.47\%), its proven pre-training alignment with Sanskrit, and its performance parity with the optimized Phase 2 models, we select \textbf{MuRIL} as the primary inference engine for our large-scale literary analysis. MuRIL's tight fold-level consistency, with all folds scoring within 1.7 percentage points of each other, suggests its learned representations generalize robustly to unseen Sanskrit verse. This reliability is the critical prerequisite for deploying the model on compositionally diverse, out-of-distribution texts like the \textit{\d Rgveda}.


\section{Literary and Corpus Analysis}

We apply the selected MuRIL model to conduct the first systematic, large-scale computational emotional analysis of two foundational Sanskrit texts.

\subsection{The \textit{V\=alm\=iki R\=am\=aya\d na}: Emotional Cartography Across All Seven K\=a\d n\d das}

The complete \textit{V\=alm\=iki R\=am\=aya\d na}, comprising 23,402 verses spanning all seven k\=a\d n\d das, was classified using the MuRIL model (fold 1). Table~\ref{tab:kanda_dominant} presents the kanda-level rasa distribution.

\begin{table}[htbp]
\centering
\small
\caption{Dominant rasa per k\=a\d n\d da in the \textit{V\=alm\=iki R\=am\=aya\d na} (\textit{rasa\_1} predictions).}
\label{tab:kanda_dominant}
\begin{tabular}{llrr}
\toprule
K\=a\d n\d da & Dominant Rasa & \% & Verses \\
\midrule
B\=ala          & V\=ira        & 26.0\% & 2,217 \\
Ayodhy\=a       & Karu\d n\=a   & \textbf{33.3\%} & 4,262 \\
Ara\d nya        & V\=ira        & 20.8\% & 2,439 \\
Ki\d skindh\=a   & V\=ira        & 30.3\% & 2,445 \\
Sundara         & V\=ira        & 27.7\% & 2,772 \\
Yuddha          & V\=ira        & \textbf{39.0\%} & 5,693 \\
Uttara          & V\=ira        & 25.5\% & 3,574 \\
\midrule
\textbf{Total}  & V\=ira        & 27.67\% & 23,402 \\
\bottomrule
\end{tabular}
\end{table}

\begin{table*}[htbp]
\centering
\normalsize
\caption{Rasa distribution across k\=a\d n\d das (verse counts).}
\label{tab:kanda_full}
\begin{tabular}{@{}lrrrrrrrrr@{}}
\toprule
K\=a\d n\d da & V & K & Sh & R & A & B & Sa & Bi & H \\
\midrule
B\=ala     & 576 & 205 & 257 & 181 & 442 & 63  & 407 & 39  & 47  \\
Ayodhy\=a  & 751 &\textbf{1419}& 631 & 164 & 417 & 129 & 551 & 108 & 92  \\
Ara\d nya   & 507 & 439 & 340 & 445 & 241 & 173 & 168 & 87  & 39  \\
Ki\d skindh\=a& 741 & 414 & 294 & 213 & 420 & 143 & 135 & 46  & 39  \\
Sundara    & 768 & 484 & 410 & 279 & 445 & 130 & 109 & 89  & 58  \\
Yuddha     &\textbf{2221}& 524 & 335 &\textbf{1278}& 659 & 309 & 213 & 94  & 60  \\
Uttara     & 911 & 463 & 405 & 439 & 535 & 158 & 504 & 69  & 90  \\
\bottomrule
\end{tabular}
\end{table*}

\textbf{Ayodhy\=a K\=a\d n\d da: the grief canon.} This k\=a\d n\d da is the sole kanda dominated by \textit{Karu\d n\=a} (33.3\%, 1,419 verses), and its dominance margin of 751 \textit{V\=ira} versus 1,419 \textit{Karu\d n\=a} is the most decisive affective skew in the entire epic. The result aligns precisely with the classical characterisation of this book as \textit{karu\d na-pradh\=ana}: it narrates the preparations for Rama's exile, Kaikey\=i's boons, Dasharatha's death, and Bharata's anguished discovery. The model's identification captures grief not merely in overtly lamentation passages but also in the resignation and restraint that saturate the dialogue. The elevated \textit{Ś\=anta} score (551 verses, 12.9\%) also reflects the stoic philosophical acceptance expressed by Rama and Sita throughout this book.

\textbf{Yuddha K\=a\d n\d da: heroism at its apex.} The largest k\=a\d n\d da (5,693 verses) shows the highest \textit{V\=ira} concentration in the entire corpus (39.0\%, 2,221 verses), confirming its role as the epic's martial climax. Critically, it is also the only k\=a\d n\d da where \textit{Raudra} rises to near-\textit{V\=ira} levels: 1,278 verses (22.4\%). This emergent \textit{V\=ira-Raudra} dyad is consistent with the aesthetic theory of \textit{rasa-sa\.mghata} (rasa confluence): in climactic battle, the hero's righteous valor (\textit{V\=ira}) and furious destruction (\textit{Raudra}) co-exist and amplify each other. The 309 \textit{Bhay\=anaka} verses (5.4\%) capture the terror experienced by Ravana's forces under Rama's assault at peak confidence scores above 95\%.

\textbf{B\=ala K\=a\d n\d da: origin stories and the cosmological register.} Beyond the expected \textit{V\=ira} dominance (26.0\%), the B\=ala K\=a\d n\d da shows a surprisingly high \textit{Ś\=anta} count (407 verses, 18.4\%), the highest of any kanda in absolute proportion. This reflects the book's cosmological prologue, creation narratives, and the serene philosophical dialogues between sages. The elevated \textit{Adbhuta} (442 verses, 19.9\%) captures the marvel of Dasharatha's yajna, the divine gift of the payasa, and the birth of Rama. This dual \textit{Ś\=anta-Adbhuta} texture is the distinctive affective signature of origin mythology.

\textbf{Ara\d nya K\=a\d n\d da: tension and dread in the forest.} Though \textit{V\=ira} nominally dominates (20.8\%), the Ara\d nya K\=a\d n\d da has the most balanced rasa distribution of any book and shows the highest \textit{Raudra} count (445 verses, 18.2\%) outside the Yuddha K\=a\d n\d da. The slaying of demons including T\=ataka, Kh\=ara, and D\=u\d sa\d na drives this intensity. The 173 \textit{Bhay\=anaka} verses (7.1\%) reflect the forests beyond Ja\=nasth\=ana; the abduction of Sita lifts \textit{Karu\d n\=a} to 439 verses (18.0\%), creating an affective swell that anticipates the grief of the Ayodhy\=a K\=a\d n\d da.

\textbf{Sundara K\=a\d n\d da: reconnaissance and longing.} Named for its beauty, the Sundara K\=a\d n\d da carries the highest per-verse \textit{Ś\d r\d ng\=ara} count proportionally (410 verses, 14.8\%) of any kanda outside Ayodhy\=a, capturing Sita's memories of Rama, Hanuman's descriptions of Lanka's grandeur, and the pathos of separation. The 484 \textit{Karu\d n\=a} verses (17.5\%) reflect Sita's captivity and anguish. The \textit{Adbhuta} score (445 verses, 16.1\%) is driven by Hanuman's ocean crossing (the top-confidence Adbhuta verses in the corpus, reaching 91.37\%), Lanka's architecture, and the \textit{Chudamani} reunion scene.

\textbf{Ki\d skindh\=a K\=a\d n\d da: alliance and the monkey kingdom.} The \textit{V\=ira} dominance (30.3\%) here is expected, but the high \textit{Adbhuta} (420 verses, 17.2\%) stands out: it arises from the marvels of the vanaras' supernatural capabilities, Sugriva's descriptions of mountain ranges, and Rama's mythological feats in defeating Vali. This kanda thus combines martial heroism with geographic and cosmological wonder.

\textbf{Uttara K\=a\d n\d da: aftermath and eternity.} The Uttara K\=a\d n\d da, the epilogue, shows the most diffuse distribution, with \textit{V\=ira} at 25.5\% but with pronounced \textit{Ś\=anta} (504 verses, 14.1\%) and substantial \textit{Karu\d n\=a} (463 verses, 13.0\%). The \textit{Ś\=anta} here is the highest count after the B\=ala K\=a\d n\d da. The \textit{Ś\=anta} reflects the philosophical discussions, Rama's return to dharmic equilibrium, and the ultimate dissolution of the narrative. The interweaving of \textit{V\=ira} and \textit{Ś\=anta} here mirrors the classical idea of \textit{v\=ira} shading into \textit{ś\=anta} at the culmination of heroic biography.

\textbf{Corpus-level rasa overview.}
Across all 23,402 verses: \textit{V\=ira} (27.67\%), \textit{Karu\d n\=a} (16.87\%), \textit{Adbhuta} (13.50\%), \textit{Raudra} (12.82\%), \textit{Ś\d r\d ng\=ara} (11.42\%), \textit{Ś\=anta} (8.92\%), \textit{Bhay\=anaka} (4.72\%), \textit{B\=ibhatsa} (2.27\%), \textit{H\=asya} (1.82\%). The low \textit{H\=asya} score (1.82\%) is consistent with the R\=am\=aya\d na's character as a \textit{mah\=ak\=avya} rather than a comic genre text; the \textit{H\=asya} that does appear is concentrated in playful exchanges (the topmost-confidence H\=asya verse reaches 98.04\% confidence) in brief dialogic interludes.

\begin{figure*}[!t]
\centering
\includegraphics[width=\textwidth]{images/Unknown-4.png}
\caption{Normalised rasa distribution across all seven k\=a\d n\d das of the \textit{V\=alm\=iki R\=am\=aya\d na}. The Ayodhy\=a K\=a\d n\d da's \textit{Karu\d n\=a} dominance and the Yuddha K\=a\d n\d da's \textit{V\=ira-Raudra} peak are clearly visible.}
\label{fig:kanda_stacked}
\end{figure*}

\subsection{Multi-Rasa Secondary Patterns in the R\=am\=aya\d na}

Analysis of the \textit{rasa\_2} (secondary) predictions reveals systematic affective undercurrents that a single-label approach would obscure:

\textbf{The V\=ira-H\=asya bond.} Almost all top-confidence \textit{V\=ira} verses (confidence above 76\%) carry \textit{H\=asya} as the secondary label at approximately 6--7\% confidence. This is not a classification error; it reflects a genuine aesthetic phenomenon in Sanskrit heroic verse. The spirited exuberance of martial description, the playful competitiveness of warriors, and the lightness of heroic speech frequently carry an undertone of spirited laughter. Sanskrit critic Abhinavagupta notes that \textit{h\=asya} naturally arises as a trace emotion (\textit{sth\=ay\=ibh\=ava}) in settings of confident power.

\textbf{The Karu\d n\=a-Ś\=anta resonance.} The top bereavement verses consistently show \textit{Ś\=anta} as secondary (5--7\% range), reflecting the classical teaching that \textit{karu\d n\=a} in its highest form resolves into the detachment of \textit{ś\=anta}. Verses describing Rama's equanimity in exile or Sita's meditative endurance in Lanka exemplify this pattern.

\textbf{Bhay\=anaka with minimal secondary ambiguity.} Top-confidence \textit{Bhay\=anaka} verses reach above 95\% primary confidence with less than 2\% on any secondary rasa, representing a remarkable degree of categorical certainty. This reflects that Sanskrit's vocabulary of terror (lexemes around \textit{bhaya}, \textit{tr\=asa}, \textit{dagdha}) is essentially non-overlapping with other registers.

\textbf{The Ś\=anta signature.} The highest-confidence classifications overall belong to \textit{Ś\=anta}, reaching 97.22\% on ritual and philosophical verses from the B\=ala K\=a\d n\d da. This indicates that the language of renunciation, sacrifice, and brahminical ethics (\textit{dharma}, \textit{yaj\~na}, \textit{k\d sam\=a}) is the most linguistically distinct rasa register, a finding consistent with the unique metrical and formulaic patterns of sacred verse.

\subsection{The \textit{\d Rgveda}, Mandala 1: A Literature of Wonder}

We performed zero-shot inference on all 1,758 hymns of \d Rgveda Mandala 1 to examine its distinctive aesthetic profile.

\begin{table}[htbp]
\centering
\small
\caption{Rasa distribution in the \textit{\d Rgveda} Mandala 1.}
\label{tab:rigveda_dist}
\begin{tabular}{lrr}
\toprule
Rasa & Count & \% \\
\midrule
Adbhuta   & 418 & \textbf{23.78\%} \\
V\=ira     & 313 & 17.80\% \\
Ś\d r\d ng\=ara & 266 & 15.13\% \\
Karu\d n\=a & 230 & 13.08\% \\
Ś\=anta   & 213 & 12.12\% \\
Raudra    & 126 & 7.17\%  \\
B\=ibhatsa &  94 & 5.35\%  \\
H\=asya   &  50 & 2.84\%  \\
Bhay\=anaka & 48 & 2.73\%  \\
\midrule
\textbf{Total} & \textbf{1,758} & \textbf{100.0\%} \\
\bottomrule
\end{tabular}
\end{table}

The \d Rgveda Mandala 1 is overwhelmingly \textit{Adbhuta}-dominant (23.78\%), driven by the cosmic marveling that pervades hymns to Indra, Agni, Varuna, and the Aśvins. Vedic poetry is fundamentally a literature of wonder: deities are praised for inconceivable feats, cosmic patterns (the sun's daily crossing, the waters' eternal flow) are invoked with astonishment, and the sacrifice itself is framed as a miraculous act of participation in divine order.

\textbf{Beauty-in-Invocation (\textit{Ś\d r\d ng\=ara}, 15.13\%).} This finding reflects the unique Vedic register of \textit{Ś\d r\d ng\=ara}, which arises not from erotic narrative but from beauty-in-invocation, seen in the radiance of U\d sa\d s (divine dawn), the elegance of soma preparation hymns, and the aesthetic pleasure of the priest-god relationship (\textit{sakhya}). The top-confidence \textit{Ś\d r\d ng\=ara} hymns in the corpus (82 to 83\% confidence) center on Dawn and on the Aśvins as beneficent beauty-bearers.

\textbf{The Meditative Mode (\textit{Ś\=anta}, 12.12\%).} The meditative, equilibrium-seeking quality of the \d Rgveda's cosmological hymns, particularly those invoking \d Rta (cosmic order) and Varu\d na's moral superintendence, maps naturally onto the serenity register. The top-confidence \textit{Ś\=anta} hymns (reaching 96.57\%) invoke Rudra for peace (\textit{ś\=antama}), appeal to cosmic law, and express the priest's serene alignment with divine will.

\textbf{Grotesque Ritual Imagery (\textit{B\=ibhatsa}, 5.35\%).} This finding reflects a genuine feature of Vedic ritual poetry: hymns dealing with demonology, sacrificial anatomy, and apotropaic (protective) magic employ visceral grotesque imagery that the model correctly identifies with the disgust register. Hymns describing the physical details of animal sacrifice, demonic anatomy, or the smashing of enemies carry imagery that Sanskrit criticism associates with \textit{v\=ibhatsa}.

\textbf{Mythological Valor (\textit{V\=ira}, 17.80\%).} The \d Rgveda's heroism is mythological and cosmic, often narrated in the perfective past tense from a worshipper's perspective. The model captures this martial intensity as a secondary but foundational mode, particularly in the Indra hymns, where the slaying of Vritra provides the primary heroic signal.

\textbf{Dread and Fury (\textit{Bhay\=anaka} and \textit{Raudra}).} The fearsome aspects of Indra and Rudra are invoked in praise rather than experienced through narrative suffering, leading to a more abstract representation of terror and fury compared to epic literature. The model correctly identifies these as smaller but significant components of the Vedic affective universe.

\begin{figure*}[!t]
\centering
\includegraphics[width=\textwidth]{images/Unknown-5.png}
\caption{Rasa distribution in \d Rgveda Mandala 1. The \textit{Adbhuta} primacy (23.78\%) constitutes a clear computational signature of the Vedic hymnal register, distinguishing it from the heroic dominance typical of Sanskrit epic narrative.}
\label{fig:genre_dist}
\end{figure*}

\section{Discussion}

The prevailing performance of MuRIL among the tested architectures is by no means incidental. Sanskrit is a morphologically intricate language characterized by an immense combinatorial space for inflected forms; its extensive use of compounds, complex \textit{sandhi} rules, and dependent case structures find few parallels in modern European languages. By pre-training on multilingual corpora that explicitly included Sanskrit, MuRIL developed underlying representations capable of subsuming these sub-morphological, language-specific patterns. This phenomenon aligns with established findings on the efficacy of domain-adaptive pre-training (Gururangan et al., 2020). Concurrently, IndicBERT's dramatic resurgence in Phase 2 highlights a critical caveat: while representational capacity for Indic scripts is a baseline requirement, it relies heavily on meticulous configuration. Appropriate loss functions (such as Focal Loss to manage class imbalance and label smoothing to dampen annotation noise) combined with finely calibrated learning rates are equally indispensable for unlocking a model's full potential in this domain.

Our LLM-ensemble annotation strategy also warrants reflection. While recent studies have demonstrated that general-purpose LLMs can effectively solve a wide range of NLP tasks in zero-shot settings (Qin et al., 2023), they have also exposed the fragility of these models when confronted with culturally dense, domain-specific tasks. Rasa classification falls squarely into this vulnerable category. To counteract this, we instituted a three-LLM consensus threshold: unanimous agreement for majority rasas and a two-of-three majority for minority rasas. This principled, democratic voting mechanism effectively dilutes individual model biases. The subsequent expert human review then acts as a final filter, eliminating systematic errors that all three models might theoretically share. The result is a corpus of a caliber that would be extraordinarily difficult, if not impossible, to assemble via crowd-sourcing for a task demanding specialized philological training.

Finally, our literary findings carry substantial implications for the broader field of digital humanities. The pronounced \textit{Karu\d n\=a} dominance identified in the Ayodhy\=a K\=a\d n\d da provides empirical, quantitative backing for traditional scholarly assessments from India and Nepal. Similarly, the extraction of an \textit{Adbhuta}-primary signature for the \d Rgveda establishes a novel computational aesthetic footprint for Vedic hymnal literature. Extending this methodology to other canonical texts, such as the Mah\=abh\=arata, K\=alid\=asa's dramas, or Bh\=aravi's \textit{K\=ir\=at\=arjun\=iya}, promises to generate a comprehensive, comparative affective atlas mapping the entirety of the classical Sanskrit tradition.

\section{Conclusion}

In this work, we introduced \texttt{SanskritRasaBank}, a pioneering corpus comprising 17,462 expert-validated Sanskrit verse annotations across nine classical \textit{rasa} categories. We compiled this dataset using a rigorous hybrid pipeline combining LLM-ensemble consensus filtering with meticulous human-expert validation drawn from eight foundational Sanskrit sources spanning the literary heritage of India and Nepal. Our systematic benchmarking reveals that MuRIL, bolstered by its dedicated pre-training on the languages of India and Nepal, delivers exceptional stability and accuracy, achieving 80.48\% (3-fold) in Phase 1 and 80.65\% in Phase 2, consistently setting the standard against all evaluated alternatives. Furthermore, our targeted hyperparameter optimization propelled IndicBERT V2 from 65.21\% to an impressive 81.45\%, a surge of over 16 percentage points, underscoring that meticulous, domain-specific tuning is just as critical as the underlying architecture for low-resource classical languages.

By deploying our highest-performing model at scale, we produced the first computational emotional cartography of both the complete \textit{V\=alm\=iki R\=am\=aya\d na} and \d Rgveda Mandala 1. The resulting per-kanda and genre-specific insights not only reaffirm centuries of traditional literary scholarship but also chart new empirical pathways for quantitative literary history. To foster continued exploration of Sanskrit's profound affective heritage, we have released all associated data, models, and code to the public domain.

\section{Limitations and Future Work}

While our annotation pipeline is explicitly designed for rigor, it retains certain irreducible limitations. Large language models may exhibit systemic blind spots when evaluating culturally intricate concepts like \textit{dhvani} (suggested meaning), which rely heavily on implicit cultural knowledge outside standard training distributions. Although our expert audit mitigates this risk, it is ultimately bounded by human bandwidth. Furthermore, the current framework operates as a single-label classification model at the primary level. Because classical rasa theory inherently recognizes verses as multi-layered phenomena (\textit{rasa-dhvani}), transitioning to a full multi-label architecture represents a vital methodological next step. 

Looking ahead, we intend to expand \texttt{SanskritRasaBank} to encompass the Mah\=abh\=arata and to explore context-aware models capable of pooling contiguous verse representations. We will also investigate whether supplementary pre-training on established digitized Sanskrit corpora (such as the DCS) can further elevate specialized model performance. A major priority for our upcoming research involves extending this computational methodology beyond metrical verse (\textit{padhyas}) into Sanskrit prose (\textit{gadhyas}), allowing us to determine whether the affective signatures of prose narratives mirror the distributional arcs of their poetic counterparts.

\section*{Statements and Declaration}

This research was conducted entirely without external financial support. No funding, grants, sponsorships, or any other form of monetary or institutional assistance were received or utilized in carrying out this work.

\begin{thebibliography}{99}


\bibitem{Conneau2020} A. Conneau, K. Khandelwal, N. Goyal, V. Chaudhary, G. Wenzek, F. Guzm\'an, E. Grave, M. Ott, L. Zettlemoyer, and V. Stoyanov. 2020. Unsupervised cross-lingual representation learning at scale. In \emph{Proceedings of the 58th Annual Meeting of the Association for Computational Linguistics (ACL)}.

\bibitem{DeepSeek2024} DeepSeek-AI. 2024. DeepSeek-V2: A strong, economical, and efficient Mixture-of-Experts Language Model. \emph{arXiv:2405.04434}.

\bibitem{Devlin2019} J. Devlin, M.W. Chang, K. Lee, and K. Toutanova. 2019. BERT: Pre-training of deep bidirectional transformers for language understanding. In \emph{Proceedings of NAACL-HLT}.

\bibitem{Doddapaneni2023} S. Doddapaneni, R. Aralikatte, G. Ramesh, S. Goyal, M.M. Khapra, A. Kunchukuttan, and P. Kumar. 2023. Towards leaving no Indic language behind: Building monolingual corpora, benchmark and models for Indic languages. In \emph{Proceedings of the 61st Annual Meeting of the Association for Computational Linguistics (ACL)}.

\bibitem{Gilardi2023} F. Gilardi, M. Alizadeh, and M. Kubli. 2023. ChatGPT outperforms crowd workers for text-annotation tasks. \emph{Proceedings of the National Academy of Sciences}, 120(30), e2305016120.

\bibitem{Gururangan2020} S. Gururangan, A. Marasovic, S. Swayamditta, K. Lo, I. Beltagy, D. Downey, and N.A. Smith. 2020. Don't stop pretraining: Adapt language models to domains and tasks. In \emph{Proceedings of the 58th Annual Meeting of the Association for Computational Linguistics (ACL)}, pages 8342--8360.

\bibitem{Hellwig2015} O. Hellwig. 2015. Using recurrent neural networks for joint lemmatisation and morphological tagging of classical Sanskrit. In \emph{Proceedings of the Workshop on Language Technologies for Digital Humanities and Cultural Heritage}.

\bibitem{Hellwig2010} O. Hellwig. 2010--present. DCS: Digital Corpus of Sanskrit. Available at: \url{https://www.sanskrit-linguistics.org/dcs/}.

\bibitem{HellwigNehrdich2018} O. Hellwig and S. Nehrdich. 2018. Sanskrit word segmentation using character-level recurrent and convolutional neural networks. In \emph{Proceedings of EMNLP}.

\bibitem{Honovich2023} O. Honovich, T. Scialom, O. Levy, and T. Schick. 2023. Unnatural instructions: Tuning language models with (almost) no human labor. In \emph{Proceedings of the 61st Annual Meeting of the Association for Computational Linguistics (ACL)}.

\bibitem{Hu2022} E.J. Hu, Y. Shen, P. Wallis, Z. Allen-Zhu, Y. Li, S. Wang, L. Wang, and W. Chen. 2022. LoRA: Low-rank adaptation of large language models. In \emph{Proceedings of the 10th International Conference on Learning Representations (ICLR)}.

\bibitem{Huet2003} G. Huet. 2003. A functional toolkit for morphological and phonological processing, application to a Sanskrit tagger. \emph{Journal of Functional Programming}, 15(4):573--614.

\bibitem{HuetKulkarni2009} G. Huet and A. Kulkarni. 2009. Reversible step-by-step Sandhi Analysis. In \emph{Proceedings of the Sanskrit Computational Linguistics Symposium (SKLT)}.

\bibitem{Kakwani2020} D. Kakwani, A. Kunchukuttan, S. Golla, N.C. Gokul, A. Bhatt, M.M. Khapra, and P. Kumar. 2020. IndicNLPSuite: Monolingual corpora, evaluation benchmarks and pre-trained multilingual language models for Indian languages. In \emph{Findings of the Association for Computational Linguistics: EMNLP 2020}.

\bibitem{Khanuja2021} S. Khanuja, D. Bansal, S. Mehtani, S. Khosla, A. Dey, B. Gopalan, and P. Talukdar. 2021. MuRIL: Multilingual representations for Indian languages. \emph{arXiv:2103.10730}.

\bibitem{Krishna2018a} A. Krishna, S. Seshadri, R. Motlani, P. Mundkur, M. Sharma, and A. Jha. 2018. Free as in free word order: An energy based model for word segmentation and morphological tagging in Sanskrit. In \emph{Proceedings of EMNLP}.

\bibitem{Krishna2021} A. Krishna, B. Santra, A. Gupta, P. Satuluri, and P. Goyal. 2021. A graph-based framework for structured prediction tasks in Sanskrit. \emph{Computational Linguistics}, 46(4):785--845.

\bibitem{Kulkarni2010} A. Kulkarni, D. Shukl, and G.N. Jha. 2010. P\=a\d ninian grammar framework for Sanskrit. In \emph{Linguistic Issues in Language Technology}, 3(1).

\bibitem{Kumar2022} A. Kumar, P. Goyal, and V. Singh. 2022. IIT-R STSA: A dataset for Sanskrit text sentiment analysis. \emph{Applied Intelligence}, 52:4317--4332.

\bibitem{Kumar2019} S. Kumar, J.P. Singh, and S.K. Dwivedi. 2019. Rasa classification of Hindi poems using part-of-speech-based and emotional features. \emph{International Journal of Natural Language Computing (IJNLC)}, 8(4).

\bibitem{LampleConneau2019} G. Lample and A. Conneau. 2019. Cross-lingual language model pretraining. In \emph{Proceedings of NeurIPS}.

\bibitem{Lin2017} T.Y. Lin, P. Goyal, R. Girshick, K. He, and P. Doll{\'{a}}r. 2017. Focal loss for dense object detection. In \emph{Proceedings of the IEEE/CVF International Conference on Computer Vision (ICCV)}.

\bibitem{Liu2019} Y. Liu, M. Ott, N. Goyal, J. Du, M. Joshi, D. Chen, O. Levy, M. Lewis, L. Zettlemoyer, and V. Stoyanov. 2019. RoBERTa: A robustly optimized BERT pretraining approach. \emph{arXiv:1907.11692}.

\bibitem{Meta2024} Meta AI. 2024. The Llama 3 Herd of Models. \emph{arXiv:2407.21783}.

\bibitem{Mishra2018} R. Mishra, S. Gupta, and R. Navarathna. 2018. Bhaavana: An emotion-based classification system for Hindi poetry. In \emph{Proceedings of the International Conference on Natural Language Processing (ICON)}.

\bibitem{Mittal2010} V. Mittal. 2010. Automatic Sanskrit segmentiser using finite state transducers. In \emph{Proceedings of the 48th Annual Meeting of the Association for Computational Linguistics: System Demonstrations}.

\bibitem{Mohammad2022} S.M. Mohammad. 2022. Ethics sheet for automatic emotion recognition and sentiment analysis. \emph{Computational Linguistics}, 48(2):239--278.

\bibitem{Moller2023} A.G. M{\o}ller, S.-J. Paasch-Colberg, et al. 2023. Is ChatGPT a good sentiment analyzer? A preliminary study. \emph{arXiv:2304.04339}.

\bibitem{OpenAI2024} OpenAI. 2024. GPT-4o System Card. \emph{arXiv:2410.21276}.

\bibitem{Pal2023} S. Pal, S. Ghosh, and A. Ekbal. 2023. SanBERTa: A pre-trained language model for Sanskrit NLP. In \emph{Proceedings of the 14th Conference on Language Resources and Evaluation (LREC-COLING)}.

\bibitem{Qin2023} C. Qin, A. Zhang, Z. Zhang, J. Chen, M. Yasunaga, and D. Yang. 2023. Is ChatGPT a general-purpose natural language processing task solver? In \emph{Proceedings of the 2023 Conference on Empirical Methods in Natural Language Processing (EMNLP)}, pages 1339--1384.

\bibitem{Sandhan2022} J.B. Sandhan, C. Shrimankar, A. Kale, and P. Goyal. 2022. Revisiting subword tokenization: A case study with sequence-to-sequence models for Sanskrit Sandhi splitting. In \emph{Proceedings of the 29th International Conference on Computational Linguistics (COLING)}, pages 3553--3563.


\bibitem{Sandhan2025} J. Sandhan, A. Barbadikar, M. Maity, P. Satuluri, T. Sandhan, R.M. Gupta, P. Goyal, and L.K. Behera. 2025. Aesthetics of Sanskrit poetry from the perspective of computational linguistics: A case study on Śik\d s\=a\d s\d taka. In \emph{Proceedings of the World Sanskrit Conference 2025: Computational Sanskrit and Digital Humanities}.

\bibitem{SatuluriBhatt2020} P. Satuluri and S. Bhatt. 2020. Simplified dependency parsing of Sanskrit. \emph{Computational Linguistics}, 46(2):411--450.

\bibitem{Singla2019} Y.K. Singla, D. Mahata, S. Aggarwal, A. Chugh, R. Maheshwari, and R.R. Shah. 2019. BHAAV: A text corpus for emotion analysis from Hindi stories. \emph{arXiv:1910.04073}.

\bibitem{Wang2023} Y. Wang, Y. Kordi, S. Mishra, A. Liu, N.A. Smith, D. Khashabi, and H. Hajishirzi. 2023. Self-instruct: Aligning language models with self-generated instructions. In \emph{Proceedings of the 61st Annual Meeting of the Association for Computational Linguistics (ACL)}.

\end{thebibliography}

\appendix
\section{Qualitative Repository Samples}
\label{sec:appendix}

Table~\ref{tab:rasa_samples} presents one representative verse for each of the nine rasas drawn from \emph{different source corpora} in \texttt{SanskritRasaBank}, illustrating the aesthetic breadth of the dataset across Sanskrit literature from Vedic hymns to classical drama. Sources include the \textit{V\=alm\=iki R\=am\=aya\d na}, K\=alid\=asa's \textit{Abhij\~n\=anaś\=akuntalam}, the \textit{Subh\=a\d sitaratnakosha}, the \textit{Vet\=alapañcavi\d m\'{s}ati}, and Mandala 1 of the \textit{\d Rgveda}.

\begin{table*}[t]
\centering
\small
\begin{tabular}{p{1.7cm} p{9.5cm} p{3cm}}
\toprule
\textbf{Rasa} & \textbf{Representative Verse (Sanskrit)} & \textbf{Source Corpus} \\
\midrule
\textbf{Ś\d r\d ng\=ara} & \skt{कान्तेन प्रहितो नवः प्रियसखीवर्गेण बद्धस्पृहः। चित्तेनोपहृतः स्मराय न समुत्स्रष्टुं गतः पाणिना। आमृष्टो मुहुरीक्षितो मुहुरभिघ्रातो मुहुर्...} & \textit{Subh\=a\d sitaratnakosha} \\
\textbf{H\=asya} & \skt{क्व वयं क्व परोक्षमन्मथो मृगशावैः समम् एधितो जनः। परिहासविजल्पितं सखे परमार्थेन न गृह्यतां वचः।} & \textit{Abhij\~n\=anaś\=akuntalam} \\
\textbf{Karu\d n\=a} & \skt{तत् साधुकृतसंधानं प्रतिसंहर सायकम्। आर्तत्राणाय वः शस्त्रं न प्रहर्तुम् अनागसि।} & \textit{Abhij\~n\=anaś\=akuntalam} \\
\textbf{Raudra} & \skt{ईदृशानां सहस्राणि सायुधानां नराधम। भक्षितानि मया रोषात् कालमाकाङ्क्षसे नु किम्।।7.68.6।।} & \textit{R\=am\=aya\d na}, Uttara K. \\
\textbf{V\=ira} & \skt{न हि तं परिपश्यामि यस्तरेत महार्णवम्। अन्यत्र गरुडाद्वायोरन्यत्र च हनूमतः।।6.1.3।।} & \textit{R\=am\=aya\d na}, Yuddha K. \\
\textbf{Bhay\=anaka} & \skt{तस्य स्कन्धे चिताधूमदग्धस्य क्रव्यगन्धिनः। सोऽपश्यल्लम्बमानं तं भूतस्येव शवं तरोः।} & \textit{Vet\=alapañcavi\d m\'{s}ati} \\
\textbf{B\=ibhatsa} & \skt{ताम्रं कार्ष्णायसं चैव तैक्ष्ण्यादेवाभ्यजायत।।1.37.19।। मलं तस्याभवत्तत्र त्रपु सीसकमेव च।} & \textit{R\=am\=aya\d na}, B\=ala K. \\
\textbf{Adbhuta} & \skt{उत स्म ते वनस्पते वातो वि वात्यग्रमित्। अथो इन्द्राय पातवे सुनु सोममुलूखल॥६॥} & \textit{\d Rgveda} Mandala 1 \\
\textbf{Ś\=anta} & \skt{शान्तम् इदम् आश्रमपदं स्फुरति च बाहुः। कुतः फलम् इहास्य अथवा भवितव्यानां द्वाराणि भवन्ति सर्वत्र।} & \textit{Abhij\~n\=anaś\=akuntalam} \\
\bottomrule
\end{tabular}
\caption{Representative samples of all nine rasas from \texttt{SanskritRasaBank}, drawn from five distinct source corpora: \textit{Vālmīki Rāmāyaṇa}, Kālidāsa's \textit{Abhijñānaśākuntalam}, the \textit{Subhāṣitaratnakoṣa}, the \textit{Vetālapañcaviṃśati}, and the \textit{Ṛgveda}. Verse texts are sourced directly from the \texttt{MERGED\_FINAL} training dataset.}
\label{tab:rasa_samples}
\end{table*}


\section{Top-Confidence Inference Outputs: \textit{V\=alm\=iki R\=am\=aya\d na}}
\label{sec:ramayana_top}

Table~\ref{tab:ramayana_top} lists the highest-confidence MuRIL prediction for each rasa class from the complete R\=am\=aya\d na inference run (23,402 verses). Confidence is the softmax probability for \textit{rasa\_1}; the secondary rasa (\textit{rasa\_2}) and its probability are shown to illustrate the affective grain of each prediction.

\begin{table*}[htbp]
\centering
\small
\caption{Highest-confidence MuRIL inference output per rasa class: \textit{V\=alm\=iki R\=am\=aya\d na}.}
\label{tab:ramayana_top}
\begin{tabular}{lrp{8.5cm}lr}
\toprule
Rasa & Conf. & Verse (excerpt) & Rasa$_2$ & Conf.$_2$ \\
\midrule
Ś\=anta    & 97.22\% & \skt{सरय्वाश्चोत्तरे तीरे यज्ञभूमिर्विधीयताम्। शान्तयश्चाभिवर्धन्तां यथाकल्पं यथाविधि।।1.8.15।।} & H\=asya & 0.70\% \\
H\=asya    & 98.04\% & \skt{तस्य तद्वचनं श्रुत्वा वायोरक्लिष्टकर्मणः। अपहास्य ततो वाक्यं कन्याशतमथाब्रवीत्।।1.32.18।।} & Ś\=anta & 0.38\% \\
B\=ibhatsa & 96.83\% & \skt{नाकुण्डली नामकुटी नास्रग्वी नाल्पभोगवान्। नामृष्टो नानुलिप्ताङ्गो नासुगन्धश्च विद्यते।।1.6.10।।} & H\=asya & 0.75\% \\
Bhay\=anaka & 95.75\% & \skt{येन येन च गच्छन्ति राक्षसा भयकर्शिताः।।3.31.19।। तेन तेन स्म पश्यन्ति राममेवाग्रतः स्थितम्।} & H\=asya & 1.11\% \\
Adbhuta   & 91.37\% & \skt{इति कृत्वा मतिं साध्वीं समुद्र श्छन्नमम्भसि।।5.1.91।। हिरण्यनाभं मैनाकमुवाच गिरिसत्तमम्।} & H\=asya & 2.74\% \\
Ś\d r\d ng\=ara & 83.75\% & \skt{ततस्तु शिंशुपामूले जानकीं पर्युपस्थिताम्। अभिवाद्याब्रवीद्दिष्ट्या पश्यामि त्वामिहाक्षताम्।।5.56.1।।} & H\=asya & 4.51\% \\
Karu\d n\=a & 81.80\% & \skt{तदलं ते वनं गत्वा क्षमं न हि वनं तव। विमृशन्निह पश्यामि बहुदोषतरं वनम्।।2.28.25।।} & H\=asya & 4.61\% \\
Raudra    & 78.77\% & \skt{रौद्रचापप्रयुक्तां तां नीलोत्पलदलप्रभाम्। शिरसाधारयद्रामो न व्यथमभ्यपद्यत।।6.100.35।।} & H\=asya & 5.87\% \\
V\=ira     & 77.66\% & \skt{ततश्शुभमतिः प्राज्ञो भ्रातुः प्रियहिते रतः। लक्ष्मणः प्रतिसंरब्धो जगाम भवनं कपेः।।4.31.10।।} & H\=asya & 6.25\% \\
\bottomrule
\end{tabular}
\end{table*}

\noindent\textbf{Observations.} The \textit{Ś\=anta} class achieves the highest primary confidence (97.22\%) with an almost negligible secondary signal (0.70\%), confirming that ritual/philosophical Sanskrit is the most linguistically crystallised rasa register. \textit{H\=asya} reaches 98.04\% with only 0.38\% secondary probability, indicating that comic vocalisations in the R\=am\=aya\d na (direct quotes of mocking speech) are the least ambiguous among all classes. The \textit{V\=ira} class, despite being the corpus-dominant rasa, achieves the \emph{lowest} top-1 confidence (77.66\%), which is consistent with \textit{V\=ira} also appearing as the most common secondary class across nearly all other rasas, demonstrating that the heroic register permeates the epic's entire linguistic fabric.

\section{Top-Confidence Inference Outputs: \textit{\d Rgveda} Mandala 1}
\label{sec:rigveda_top}

Table~\ref{tab:rigveda_top} lists the highest-confidence MuRIL prediction per rasa for \d Rgveda Mandala 1 (1,758 hymns). Vedic accent marks were stripped during preprocessing; verses transcribed here are as they appear in the cleaned input.

\begin{table*}[htbp]
\centering
\small
\caption{Highest-confidence MuRIL inference output per rasa class: \textit{\d Rgveda} Mandala 1.}
\label{tab:rigveda_top}
\begin{tabular}{lrp{8.5cm}lr}
\toprule
Rasa & Conf. & Hymn excerpt & Rasa$_2$ & Conf.$_2$ \\
\midrule
H\=asya       & 97.65\% & \skt{वच्यन्ते वां ककुहासो जूर्णायामधि विष्टपि। यद् वां रथो विभिष्पतात्॥३॥} & Raudra & 0.40\% \\
Ś\=anta       & 96.57\% & \skt{कद् रुद्राय प्रचेतसे मीह्ळुष्टमाय तव्यसे। वोचेम शंतमं हृदे॥१॥} & H\=asya & 0.84\% \\
B\=ibhatsa    & 96.49\% & \skt{अंसेष्वा वः प्रपथेषु खादयो ऽक्षो वश्चक्रा समया वि वावृते॥९॥} & H\=asya & 0.92\% \\
Bhay\=anaka & 95.50\% & \skt{अभिव्रजन्नक्षितं पाजसा रजः स्थातुश्चरथं भयते पतत्रिणः॥५॥} & H\=asya & 1.19\% \\
Adbhuta      & 90.76\% & \skt{उत स्म ते वनस्पते वातो वि वात्यग्रमित्। अथो इन्द्राय पातवे सुनु सोममुलूखल॥६॥} & H\=asya & 3.08\% \\
Ś\d r\d ng\=ara & 82.81\% & \skt{या दस्रा सिन्धुमातरा मनोतरा रयीणाम्। धिया देवा वसुविदा॥२॥} & H\=asya & 4.38\% \\
Karu\d n\=a  & 79.45\% & \skt{शृण्वन्ता वामवसे जोहवीमि वृधे च नो भवतं वाजसातौ॥१२॥} & H\=asya & 8.01\% \\
Raudra       & 76.34\% & \skt{उग्रो ययिं निरपः स्रोतसासृजद्वि शुष्णस्य दृंहिता ऐरयत्पुरः॥११॥} & H\=asya & 6.98\% \\
V\=ira        & 75.86\% & \skt{तुञ्जेतुञ्जे य उत्तरे स्तोमा इन्द्रस्य वज्रिणः। न विन्धे अस्य सुष्टुतिम्॥७॥} & H\=asya & 6.75\% \\
\bottomrule
\end{tabular}
\end{table*}

\noindent\textbf{Observations.} The \d Rgveda pattern closely mirrors the R\=am\=aya\d na in that \textit{Ś\=anta} and \textit{H\=asya} achieve the highest single-class confidence, while \textit{V\=ira} is again lowest, which is consistent with the cross-genre finding that the heroic register in Sanskrit is linguistically diffuse. Notably, the top \textit{Ś\=anta} hymn invokes Rudra for peace (\textit{ś\=antama}), a direct instantiation of the pacificatory hymnal mode unique to Vedic literature. The top-confidence \textit{Bhay\=anaka} hymn describes a cosmic chariot causing universal fear (\textit{bhayate patatriṇaḥ}), and the top \textit{Adbhuta} verse invokes the Soma-grinding mortar as a cosmically wondrous object. Both of these reflect the distinctive register of Vedic ritual marveling that differentiates this corpus from epic narrative.

\end{document}